\section{绪论}

\subsection{选题背景与意义}
本次综合设计选题为``基于 STM32 微控制器的智能系统设计''方向，
具体实现融合了备选题目``个人语音助手''与``智能运动检测''两案要求，
以\textbf{智能手表原型}为形态组织在一块开发板上：
以语音对话作为人机交互入口，以运动计步、跌倒粗判与环境温度作为持续感知，
以彩屏表盘作为统一呈现界面。所谓``原型''，指系统当前运行在 STM32 开发板与外接模组上，
以最少的自制硬件完成完整可演示的功能闭环，验证的是架构与软件方案，而非可穿戴产品形态。

设计覆盖指导书要求的完整链条——信息采集（麦克风声音、六轴姿态、体温级温度、按键）、
信息处理与控制（STM32F407 上的计步算法、跌倒判定、语音组包与本地指令执行）、
数据传输与显示（UART 自定义帧、Wi-Fi/HTTPS 上云、LVGL 表盘）——
并在此基础上引入``云--边--端''三层协同：重负载的语音识别模型运行在云端，
嵌入式端只保留确定能在片内资源内完成、且断网也必须可用的功能。

\subsection{方案选择}
语音识别有两条技术路线：端侧小模型（关键词识别，离线可用）与云端大模型（自由问答，需联网）。
综合考虑课设周期、F407 的 168\,MHz/192\,KB RAM 算力与演示效果，
本设计采用\textbf{云边协同}路线：F407 只负责采集与可靠传输，
ESP32-C3 负责联网、TLS 与协议转换，模型推理交给云端现成的语音大模型服务。
这样端侧软件保持确定性，而交互能力上限由云端模型决定，
同时通过``断网降级''设计保证本地基础功能（走时、计步、测温）的网络无关性，
恰好形成一组可在答辩中现场演示的对照实验。

\subsection{主要工作与报告结构}
本人完成的主要工作包括：系统总体方案与云边端划分；F407 侧 FreeRTOS
任务划分与各外设驱动整合（FSMC 屏、外部 SRAM、ADC/DMA 录音、I$^2$C、单总线、RTC、WS2812 点阵）；
LVGL 六行主页与各子页，以及中文字体；F407 与 C3 之间的自定义可靠帧协议；
C3 侧语音流式上传与天气服务接入；以及覆盖主机仿真测试、烧录校验、实板联调的验证流程。

报告第 2 章为系统总体设计，第 3 章为硬件设计，第 4 章为软件与核心算法实现，
第 5 章为系统测试与效果分析，第 6 章总结与展望；
附录摘录三段开发对话，并给出核心代码节选。
